\section{Executive Summary}

The INTISAT Image Sensor Microscopy (ISM) payload is a CubeSat-class in-situ computational microscope implementing Fourier Ptychographic Microscopy (FPM), a computational imaging technique that reconstructs a high-resolution, complex-valued image (amplitude plus phase) from a sequence of lower-resolution intensity measurements taken under angularly diverse illumination.

The concept is validated at bench level. The software pipeline (capture $\rightarrow$ FPM reconstruction $\rightarrow$ super-resolution via RealESRGAN) runs on a Raspberry Pi 4 and produces demonstrably enhanced images. The repository contains approximately 18~GB of real experimental data, production-grade Python code, and extensive sensor/algorithm documentation. No hardware design files (PCB, schematic, or Gerber) exist in the repository.

This document closes the gap between the bench-level demonstrator and a flight-credible CubeSat payload by:
\begin{itemize}[leftmargin=1.5em]
    \item defining the hardware architecture and board partitioning,
    \item providing mission-level parameter tables and budgets,
    \item proposing an onboard versus ground processing split,
    \item identifying the differentiators with respect to prior work, and
    \item specifying the validation and test plan required for credibility.
\end{itemize}

\textbf{Recommended architecture:} Raspberry Pi Compute Module 4 (CM4) on a custom carrier board, with a separate sensor board and a microcontroller supervisor (Option D).
